\begin{theorem}[Интегрируемость сложной функции]
	Если $g\in\mc R[a,b],\ (\forall x\in[a,b])\ m\le g(x)\le M,\ f\text{ непрерывна на }[m,M]$, то сложная функция $h=f\circ g\in\mc R[a,b]$.
\end{theorem}
\begin{proof}
	$f$ непрерывна на $[m,M]\Ra f$ равномерно непрерывна на $[m,M]$.\\
	Также введём $K=\smo{\sup_{\sus[m,M]}} f(y)$, далее пригодится.
	\[(\forall\eps>0)(\ex\delta>0,\delta<\eps)(\fa s,t\in[m,M],|s-t|<\delta)\ |f(s)-f(t)|<\eps\]
	По критерию интегрируемости: $(\ex P)\ U(P,f)-L(P,f)<\delta^2$\\
	Ещё 
	\[U(P,h)-L(P,h)=\sum_{i=1}^n(M_i(h)-m_i(h))\tr x_i=(\sum_{i\in A}+\sum_{i\in B})(M_i(h)-m_i(h))\tr x_i\]
	\[i\in A\lra M_i(g)-m_i(g)<\delta,\ B=\{1,\dots,n\}\setminus A\]
	\[\sum_{i\in A}(M_i(h)-m_i(h))\tr x_i\le\eps\cdot(b-a)\]
	\begin{multline*}
		\delta\cdot\sum_{i\in B}\tr x_i\le\sum_{i\in B}(M_i(g)-m_i(g))\tr x_i\le U(P,g)-L(P,g)<\delta^2\Ra\\
		\Ra\sum_{i\in B}\tr x_i<\delta\Ra\sum_{i\in B}(M_i(h)-m_i(h))\tr x_i\le 2K\sum_{i\in B}<2K\cdot\eps
	\end{multline*}
	Получили ограничение на разность сумму Дарбу, а значит и получили требуемое.
\end{proof}
\begin{corollary}
	\begin{enumerate}
		\item $f_1,f2\in\mc R[a,b]\Ra f_1\cdot f_2\in\mc R[a,b]$\\
		\item $\displaystyle f_1\in\mc R[a,b]\Ra|f(x)|^2\in\mc R[a,b],\ \bigg|\int_a^bf(x)dx\bigg|\le\int_a^b\bigg|f(x)\bigg|dx$
	\end{enumerate}
\end{corollary}
\begin{proof}
	\begin{enumerate}
		\item $g\in\mc R[a,b],\ f(y)=y^2\Ra g^2\in\mc R[a,b]$
		\[(f_1\cdot f_2)(x)=\frac12((f_1+f_2)^2-f_1^2-f_2^2)(x)\]
		Получили представление $(f_1\cdot f_2)$ в виде интегрируемых функций.
		\item Модуль - непрерывная функция $\Ra|f|\in\mc R[a,b],\ c=\sign\int_a^bf(x)dx$
		\[\bigg|\int_a^bf(x)dx\bigg|=c\int_a^bf(x)dx=\int_a^bc\cdot f(x)dx\le\int_a^b\bigg|f(x)\bigg|dx\]
	\end{enumerate}
	Доказали требуемое.
\end{proof}
\begin{example}
	\[f(x)=
	\begin{cases}
		1,\ x\in\Q\\
		0,\ x\notin\Q
	\end{cases}
	\]
\end{example}

\subsection{Другое определение интеграла Римана}
\begin{definition}
	\textit{Интегральной суммой (Римана)} называется:
	\[S(P,f,\{t_i\})=\sum_{i=1}^nf(t_i)\cdot\tr x_i,\text{ где }P: a=x_0<\dots<x_n=b,\ t_i\in[x_{i-1},x_i],\ i=1,\dots,n\]
\end{definition}
\begin{definition}
	Число $A\in\R$ называется \textit{пределом интегральных сумм}, если:
	\[A=\lim_{\tr(P)\ra0}S(P,f,\{t_i\}), \text{ если }(\fa\eps>0)(\ex\delta>0)(\fa P,\ \tr(P)<\delta)(\fa\{t_i\},\ t_i\in[x_{i-1},x_i])\]
\end{definition}
\begin{lemma}
	Если $\ex$ предел интегральных сумм $S(P,f,\{t_i\})$, то $f$ ограничена.
\end{lemma}
\begin{proof}
	От противного. Пусть $f$, для определённости, неограничена сверху.
	\[(\fa A\in\R)(\ex\eps>0)(\fa\delta>0)(\ex P,\tr(P)<\delta)(\ex\{t_i\},\ t_i\in[x_{i-1},x_i])\ |A-S(P,f,\{t_i\})|\ge\eps\]
	\[S(P,f,\{t_i\})>n,\ \tr(P)<\frac1n,\ \ex\ i_0 \text{ такое, что }f\text{ неограничена на }[x_{i_0-1},x_{i_0}]\]
	\[t_i\in[t_{i-1},t_i],\ i\neq i_0,\ \v S(P,f,\{t_i\})=\sum_{i\neq i_0}f(t_i)\cdot\tr x_i\]
	\[\ex t_{i_0}\in[x_{i_0-1},x_i]:f(t_{i_0})>\frac{n-\v S(P,f,\{t_i\})}{\tr x_{i_0}}\]
	\[s(P,f,\{t_i\})=\sum_{i=1}^nf(t_i)\cdot\tr x_i=f(t_{i_0})\cdot\tr x_i+\v S(P,f,\{t_i\})>n\]
	Получили отрицание существования предела, значит лемма доказана.
\end{proof}
\begin{theorem}[Интеграл как предел интегральных сумм]\label{8_th1}
	$f\in\mc R[a,b]$ тогда и только тогда, когда $\ex$ предел интегральных сумм, причём:
	\[\int_a^bf(x)dx=\lim_{\tr(p)\ra0}S(P,f,\{t_i\})=A\]
\end{theorem}
\begin{proof}
	Достаточность:
	\[(\fa\eps>0)(\ex\delta>0)(\fa P,\ \tr(P)<\delta)(\fa\{t_i\},\ t_i\in[x_{i-1},x_i])\]
	Зафиксируем $\eps$, получим по нему $\delta$, зафиксируем $P$.
	\[A-\frac\eps3<\sum_{i=1}^nf(t_i)\tr x_i<A+\frac\eps3\Ra U(P,f)-L(P,f)\le \frac{2\eps}{3}<\eps\Ra f\in\mc R[a,b]\]
	\[A-\frac\eps3\le L(P,f)\le\int_a^bf(x)dx\le U(P,f)\le A+\frac\eps3\Ra\bigg|A-\int_a^bf(x)dx\bigg|\le\frac\eps3\]
	Достаточность доказана.\\\\
	Необходимость:\\
	Пусть $f(x)$ неотрицательная на $[a,b]$.\\
	Хотим доказать:
	\[(\fa\eps>0)(\ex\delta>0)(\fa P,\tr(P)<\delta)(\fa\{t_i\},\ t_i\in[x_{i-1},x_i])\bigg|\int_a^bf(x)dx-\sum_{i=1}^hf(t_i)dx_i\bigg|<\eps\]
	По критерию интегрируемости:
	\[(\ex P':a=x_0'<\dots<x_n'=b)\ U(P',f)<\int_a^bf(x)dx+\frac\eps4\]
	Положим: $M:=\smo{\sup_{[a,b]}}f(x),\ \delta_1:=\frac{\eps}{8Mn}$, возьмём $\fa P:\tr(P)<\delta_1$\\
	\[U(P,f)=\sum_{j=1}^nM_j\tr x_j=\bigg(\sum_{j:\ex k:x_k'\in[x_{j-1},x_j]}+\sum_{j:\nexists k}\bigg)M_j\tr x_j=S_1+S_2\]
	Оценим $S_1$:\\
	Заметим, что количество отрезков из разбиения $P$, на которых есть хотя бы одна точка из $P'$ явно $\le2n$, так как в $P'$ всего $n$ точек, а сама точка максимум может находиться на 2 отрезках $P$.
	\[S_1=\sum_{j:\ex k:x_k'\in[x_{j-1},x_j]}M_j\tr x_j\le M\cdot 2n\cdot \tr(P)<\frac\eps4\]
	Теперь оценим $S_2$:
	\[S_2=\sum_{j:\nexists k}\sup_{[x_{j-1},x_j]}f(x)\tr x_j\le\sum_{i=1}^nM_i'\tr x_i'<\int_a^bf(x)dx\]
	Получили:
	\[(\fa\eps>0)(\ex\delta_1>0)(\fa P,\tr(P)<\delta)\ U(P,f)<\int_a^bf(x)dx+\frac\eps2\]
	Получили часть доказательства, только для неотрицательной функции и только для верхней суммы. Теперь раскроем на остальные случаи.\\
	Если $f$ произвольного знака, то $\ex m:f_1(x)=f(x)+m\ge0$\\
	Для $f_1$ получим:
	\[U(P,f_1)<\int_a^bf_1(x)dx+\frac\eps2\]
	\[U(P,f)+m(b-a)<\int_a^bf(x)dx+m(b-a)\]
	Для произвольного знака получили такую же штуку.\\
	Теперь возьмём $f_2=-f$:
	\[(\ex\delta_2<0)(\fa P,\ \tr(P)<\delta_2)\ U(P,f_2)<\int_a^bf_2(x)dx+\frac\eps2,\ U(P,f)=-L(P,f)<-\int_a^bf(x)dx+\frac\eps2\]
	\[(\fa\eps>)(\ex\delta_2)(\fa P,\ \tr(P)<\delta_2)\ L(P,f)>\int_a^b-\frac\eps2\]
	Теперь возьмём $\delta=\min(\delta_1,\delta_2),\ \fa P,\ \tr(P)<\delta$
	\[\int_a^bf(x)dx-\frac\eps2<L(P,f)\le\sum_{i=1}^nf(t_i)dx_i\le U(P,f)<\int_a^bf(x)dx+\frac\eps2\]
	Получили требуемое.
\end{proof}
\begin{note}
	Достаточность теоремы (\ref{8_th1}) справедлива и для интеграла Римана - Стилтьеса. Необходимость верна лишь при дополнительном условии непрерывности $f$ или $\alpha$.
\end{note}